\appendix

\section{Proof Details}

\subsection{Deterministic Midpoint Localization}
\label{app:midpoint-localization}

\begin{lemma}[Uniform polynomial derivative envelope]
\label{lem:derivative-envelope}
Fix an integer $d\geq1$, a radius $R>0$, and a compact set
$\Theta\subseteq\mathbb R$, and define $\phi_\alpha$, $B$, and
$L_{d,R,\Theta}$ as defined above.  Then every
$\alpha\in[-R,R]^d$ and $\theta\in\Theta$ satisfy
\[
\left|\partial_\theta\phi_\alpha(\theta)\right|
\leq dB^{d-1}+R\sum_{k=1}^{d-1}kB^{k-1}.
\]
Consequently,
\[
L_{d,R,\Theta}
\leq dB^{d-1}+R\sum_{k=1}^{d-1}kB^{k-1}<\infty.
\]
For $d=1$, the sum is empty and $L_{1,R,\Theta}=1$.
\end{lemma}

\begin{proof}
Compactness of $\Theta$ makes
$\sup_{\theta\in\Theta}|\theta|$ finite, so $B$ is finite and at least
one.  Differentiating gives
\[
\partial_\theta\phi_\alpha(\theta)
=d\theta^{d-1}+\sum_{k=1}^{d-1}k\alpha_k\theta^{k-1}.
\]
Since $|\alpha_k|\leq R$ and $|\theta|\leq B$, the triangle inequality
gives, term by term,
\[
\begin{aligned}
\left|\partial_\theta\phi_\alpha(\theta)\right|
&\leq d|\theta|^{d-1}
  +\sum_{k=1}^{d-1}k|\alpha_k||\theta|^{k-1}\\
&\leq dB^{d-1}+R\sum_{k=1}^{d-1}kB^{k-1}.
\end{aligned}
\]
The right-hand side is a finite sum of finite quantities and is uniform
over the coefficient cube and $\Theta$.  Taking the defining supremum
proves the bound and finiteness of $L_{d,R,\Theta}$.  If $d=1$, then
$\phi_\alpha(\theta)=\theta+\alpha_0$ and
$\partial_\theta\phi_\alpha\equiv1$, proving the exact identity.
\end{proof}

\begin{proposition}[Borel root hitting for arbitrary interval endpoints]
\label{prop:borel-root-event}
Fix an integer $d\geq1$, a radius $R>0$, and a compact set
$\Theta\subseteq\mathbb R$, with $\phi_\alpha$ and $Z_\alpha$ as in the
setup.  For every $I\in\mathscr I(\Theta)$, under any open, closed, or
half-open endpoint convention,
\[
\mathcal E_I
:=\{\alpha\in[-R,R]^d:Z_\alpha\cap I\neq\varnothing\}
\]
is a Borel subset of $[-R,R]^d$.
\end{proposition}

\begin{proof}
Let $a=\inf I$ and $b=\sup I$.  Both are finite, and $a<b$.  For
$n\geq1$, set
\[
\varepsilon_n:=\frac{b-a}{2(n+1)},
\qquad
a_n:=
\begin{cases}
a,&a\in I,\\
a+\varepsilon_n,&a\notin I,
\end{cases}
\qquad
b_n:=
\begin{cases}
b,&b\in I,\\
b-\varepsilon_n,&b\notin I,
\end{cases}
\]
and $K_n=[a_n,b_n]$.  Because
$\varepsilon_n<(b-a)/2$, every $K_n$ is a nonempty compact interval
contained in $I$.  Moreover, $K_n\subseteq K_{n+1}$ and
\[
\bigcup_{n\geq1}K_n=I.
\]
An endpoint is retained in every $K_n$ exactly when it belongs to $I$,
and every interior point is eventually captured.  Thus the exhaustion
covers all four endpoint conventions.

For $\alpha\in[-R,R]^d$, define
\[
F_n(\alpha):=\min_{\theta\in K_n}|\phi_\alpha(\theta)|.
\]
The minimum exists because $K_n$ is nonempty and compact and the
polynomial is continuous.  For any
$\alpha,\gamma\in[-R,R]^d$, put
\[
\delta:=\sup_{\theta\in K_n}
|\phi_\alpha(\theta)-\phi_\gamma(\theta)|.
\]
Pointwise on $K_n$,
$|\phi_\gamma(\theta)|-\delta\leq|\phi_\alpha(\theta)|
\leq|\phi_\gamma(\theta)|+\delta$.  Taking minima and then interchanging
$\alpha$ and $\gamma$ gives
\[
\begin{aligned}
|F_n(\alpha)-F_n(\gamma)|
&\leq\sup_{\theta\in K_n}
  |\phi_\alpha(\theta)-\phi_\gamma(\theta)|\\
&\leq\sum_{k=0}^{d-1}|\alpha_k-\gamma_k|B^k.
\end{aligned}
\]
Hence $F_n$ is continuous, and
\[
\mathcal E_{I,n}:=\{\alpha\in[-R,R]^d:F_n(\alpha)=0\}
\]
is closed relative to the coefficient cube.  Since the minimum is
attained, $F_n(\alpha)=0$ exactly when $\phi_\alpha$ has a zero in $K_n$.
The compact exhaustion therefore yields
\[
\mathcal E_I=\bigcup_{n\geq1}\mathcal E_{I,n},
\]
a countable union of closed sets.  The argument only asserts existence of
a zero on each compact set and never chooses a root as a function of
$\alpha$.
\end{proof}

\begin{proposition}[Measurable midpoint intercept slab]
\label{prop:midpoint-slab}
Fix an integer $d\geq1$, a radius $R>0$, a compact set
$\Theta\subseteq\mathbb R$, and $I\in\mathscr I(\Theta)$.  Under
Lemma~\ref{lem:derivative-envelope}, define, with an empty sum when
$d=1$,
\[
c_I(\beta):=m_I^d+\sum_{k=1}^{d-1}\alpha_km_I^k,
\qquad \beta=(\alpha_1,\ldots,\alpha_{d-1}),
\]
and
\[
J_I(\beta):=
\left[-c_I(\beta)-\frac{L_{d,R,\Theta}|I|}{2},
      -c_I(\beta)+\frac{L_{d,R,\Theta}|I|}{2}\right].
\]
Then $c_I$ and both endpoints of $J_I$ are Borel measurable,
$|J_I(\beta)|=L_{d,R,\Theta}|I|$, and
\[
\{(a_0,\beta):a_0\in J_I(\beta)\}
=\left\{(a_0,\beta):
|a_0+c_I(\beta)|\leq\frac{L_{d,R,\Theta}|I|}{2}\right\}
\]
is Borel.  For every $\alpha=(\alpha_0,\beta)\in[-R,R]^d$,
\[
Z_\alpha\cap I\neq\varnothing
\quad\Longrightarrow\quad
|\phi_\alpha(m_I)|
\leq\frac{L_{d,R,\Theta}|I|}{2}
\quad\Longleftrightarrow\quad
\alpha_0\in J_I(\beta).
\]
The implication holds for every endpoint convention and requires neither
simple nor transverse roots.
\end{proposition}

\begin{proof}
Again write $a=\inf I$ and $b=\sup I$.  Since $a<b$, the midpoint
$m_I=(a+b)/2$ lies strictly between $a$ and $b$, and hence belongs to
$I$ for every endpoint convention.  If $\theta_*\in I$, the closed
segment joining $\theta_*$ and $m_I$ lies in $I\subseteq\Theta$, and
\[
|m_I-\theta_*|\leq\frac{b-a}{2}=\frac{|I|}{2}.
\]

Suppose $\theta_*\in Z_\alpha\cap I$.  If $\theta_*=m_I$, the desired
bound has left-hand side zero.  Otherwise, the one-dimensional
mean-value inequality applies on the segment: polynomiality gives
continuity and differentiability, and Lemma~\ref{lem:derivative-envelope}
bounds the derivative there.  Since $\phi_\alpha(\theta_*)=0$,
\[
\begin{aligned}
|\phi_\alpha(m_I)|
&=|\phi_\alpha(m_I)-\phi_\alpha(\theta_*)|\\
&\leq L_{d,R,\Theta}|m_I-\theta_*|\\
&\leq\frac{L_{d,R,\Theta}|I|}{2}.
\end{aligned}
\]
This retains the exact midpoint factor.

The additive-intercept form gives
\[
\phi_\alpha(m_I)=\alpha_0+c_I(\beta),
\]
so the evaluation bound is equivalent to
$\alpha_0\in J_I(\beta)$.  The endpoint difference of $J_I$ is exactly
$L_{d,R,\Theta}|I|$.  For fixed $I$, $c_I$ is a polynomial in $\beta$,
so it and the interval endpoints are continuous.  The map
$(a_0,\beta)\mapsto|a_0+c_I(\beta)|$ is continuous, and the slab is the
inverse image of the closed interval
$[0,L_{d,R,\Theta}|I|/2]$.

When $d=1$, $\beta$ is the empty tuple,
$c_I(\beta)=m_I$, $L_{1,R,\Theta}=1$, and
\[
J_I=[-m_I-|I|/2,-m_I+|I|/2].
\]
Thus the formulas remain literal on the one-point conditioning space.
If an included endpoint is a root, the distance bound may be an equality;
an excluded endpoint is absent from the event, while arbitrarily
endpoint-near roots obey the same weak inequality.  The proof never
divides by $|\partial_\theta\phi_\alpha(\theta_*)|$ and never uses root
uniqueness, so tangent and multiple roots are covered.  Together with
Proposition~\ref{prop:borel-root-event}, this supplies a Borel root event
contained in a jointly Borel intercept slab of exact width
$L_{d,R,\Theta}|I|$.
\end{proof}

\subsection{Averaged Conditional-Density Integration}
\label{app:averaged-integration}

\begin{lemma}[Conditional integration over the midpoint slab]
\label{lem:slab-disintegration}
Under Assumption~\ref{assump:averaged-intercept-density},
Proposition~\ref{prop:borel-root-event}, and
Proposition~\ref{prop:midpoint-slab}, fix
$\mu\in\mathcal D$ and $I\in\mathscr I(\Theta)$.  Then the root event
and midpoint slab are Borel, the function
\[
(a_0,\beta)\longmapsto
\mathbf1\{a_0\in J_I(\beta)\}f_\mu(a_0\mid\beta)
\]
is nonnegative and jointly measurable, and
\[
\Pr_\mu(Z_\alpha\cap I\neq\varnothing)
\leq
\int\!\int_{\mathbb R}
\mathbf1\{a_0\in J_I(\beta)\}
f_\mu(a_0\mid\beta)\,da_0\,\pi_\mu(d\beta).
\]
\end{lemma}

\begin{proof}
Proposition~\ref{prop:borel-root-event} makes
$\mathcal E_I=\{Z_\alpha\cap I\neq\varnothing\}$ measurable, while
Proposition~\ref{prop:midpoint-slab} makes
$\mathcal S_I=\{(a_0,\beta):a_0\in J_I(\beta)\}$ Borel and gives
$\mathcal E_I\subseteq\mathcal S_I$.  Hence
\[
\Pr_\mu(\mathcal E_I)\leq\Pr_\mu(\mathcal S_I).
\]
The indicator of $\mathcal S_I$ is jointly Borel.  Multiplying it by the
jointly measurable nonnegative density supplied by
Assumption~\ref{assump:averaged-intercept-density} gives the asserted
integrand.  Regular conditional-density disintegration yields
\[
\Pr_\mu(\mathcal S_I)
=\int\!\int_{\mathbb R}
\mathbf1\{a_0\in J_I(\beta)\}
f_\mu(a_0\mid\beta)\,da_0\,\pi_\mu(d\beta).
\]
The density identity is needed only for $\pi_\mu$-almost every fiber;
the remaining fibers contribute zero to the outer integral.  The
conditional intercept is supported on $[-R,R]$ and its density is
extended by zero, so the formula remains exact even if $J_I(\beta)$
extends outside the cube.  Tonelli's theorem for this nonnegative jointly
measurable integrand also makes the inner integral a measurable
extended-valued function of $\beta$.
\end{proof}

\begin{lemma}[Averaged slice-cap control]
\label{lem:averaged-slice-control}
Under Assumption~\ref{assump:averaged-intercept-density},
Proposition~\ref{prop:midpoint-slab}, and
Lemma~\ref{lem:slab-disintegration}, fix
$\mu\in\mathcal D$ and $I\in\mathscr I(\Theta)$.  For
$\pi_\mu$-almost every $\beta$,
\[
\int_{\mathbb R}
\mathbf1\{a_0\in J_I(\beta)\}
f_\mu(a_0\mid\beta)\,da_0
\leq K_\mu(\beta)L_{d,R,\Theta}|I|,
\]
and
\[
\int\!\int_{\mathbb R}
\mathbf1\{a_0\in J_I(\beta)\}
f_\mu(a_0\mid\beta)\,da_0\,\pi_\mu(d\beta)
\leq L_{d,R,\Theta}|I|
\int K_\mu(\beta)\,\pi_\mu(d\beta)<\infty.
\]
\end{lemma}

\begin{proof}
For fixed $\mu$, the definition of the class envelope gives
\[
0\leq\int K_\mu(\beta)\,\pi_\mu(d\beta)
\leq\bar\kappa_{\mathcal D}<\infty.
\]
Because $K_\mu$ is nonnegative, measurable, and possibly
extended-valued, this implies $K_\mu(\beta)<\infty$ almost surely.  To
see this explicitly, let
$N_\mu^\infty=\{\beta:K_\mu(\beta)=\infty\}$.  For every integer
$n\geq1$,
\[
\int K_\mu\,d\pi_\mu\geq n\pi_\mu(N_\mu^\infty),
\]
so finiteness forces $\pi_\mu(N_\mu^\infty)=0$.

Intersect this full-measure set with the full-measure set on which
$f_\mu(\cdot\mid\beta)$ is the conditional density and $K_\mu(\beta)$
is its Lebesgue essential supremum.  On the resulting full-measure set,
\[
f_\mu(a_0\mid\beta)\leq K_\mu(\beta)
\quad\text{for Lebesgue-almost every }a_0.
\]
Proposition~\ref{prop:midpoint-slab} gives
$|J_I(\beta)|=L_{d,R,\Theta}|I|<\infty$.  Therefore
\[
\begin{aligned}
\int_{\mathbb R}\mathbf1\{a_0\in J_I(\beta)\}
f_\mu(a_0\mid\beta)\,da_0
&=\int_{J_I(\beta)}f_\mu(a_0\mid\beta)\,da_0\\
&\leq K_\mu(\beta)|J_I(\beta)|\\
&=K_\mu(\beta)L_{d,R,\Theta}|I|.
\end{aligned}
\]
Null fibers, including a fiber assigned $K_\mu=\infty$, contribute
nothing to the outer integral.  By Lemma~\ref{lem:slab-disintegration}
and Tonelli, the left-hand inner integral is measurable and nonnegative.
Integrating the last display against the arbitrary Borel probability
$\pi_\mu$ proves the second inequality.  This uses no density of
$\pi_\mu$ and never replaces $K_\mu(\beta)$ by an essential supremum
over $\beta$.
\end{proof}

\begin{proposition}[Uniform root hitting under an averaged intercept envelope]
\label{prop:averaged-root-hitting}
Under Assumption~\ref{assump:averaged-intercept-density},
Proposition~\ref{prop:borel-root-event},
Proposition~\ref{prop:midpoint-slab},
Lemma~\ref{lem:slab-disintegration}, and
Lemma~\ref{lem:averaged-slice-control}, every
$\mu\in\mathcal D$ and $I\in\mathscr I(\Theta)$ satisfy
\[
\Pr_\mu(Z_\alpha\cap I\neq\varnothing)
\leq L_{d,R,\Theta}|I|
\int K_\mu(\beta)\,\pi_\mu(d\beta)
\leq L_{d,R,\Theta}|I|\bar\kappa_{\mathcal D}.
\]
Consequently,
\[
C_{\mathcal D}
\leq L_{d,R,\Theta}\bar\kappa_{\mathcal D}<\infty.
\]
\end{proposition}

\begin{proof}
Lemma~\ref{lem:slab-disintegration} bounds the root probability by the
nonnegative double integral over the midpoint slab, and
Lemma~\ref{lem:averaged-slice-control} bounds that integral by the first
right-hand side.  For the fixed law $\mu$,
\[
\int K_\mu(\beta)\,\pi_\mu(d\beta)
\leq\sup_{\nu\in\mathcal D}
\int K_\nu(\beta)\,\pi_\nu(d\beta)
=\bar\kappa_{\mathcal D}.
\]
Thus each fiber cap is integrated before the class supremum is applied.
Since $|I|>0$, division by $|I|$ gives
\[
\sup_{I\in\mathscr I(\Theta)}
\frac{\Pr_\mu(Z_\alpha\cap I\neq\varnothing)}{|I|}
\leq L_{d,R,\Theta}\bar\kappa_{\mathcal D}.
\]
Taking the outer supremum over $\mu\in\mathcal D$ proves the claim.
Finiteness follows from Lemma~\ref{lem:derivative-envelope} and
Assumption~\ref{assump:averaged-intercept-density}.

When $d=1$, $\beta$ is the empty tuple and $\pi_\mu$ is the point mass
on its one-point space.  Proposition~\ref{prop:midpoint-slab} gives
$L_{1,R,\Theta}=1$ and a constant interval $J_I$ of length $|I|$; the
outer integrals reduce to evaluation at one point.  Thus the argument is
the ordinary one-dimensional density bound and introduces no hidden
positive-dimensional marginal condition.  The preceding results also
show explicitly that zero extension, unbounded slice caps, singular or
discrete higher-coordinate marginals, and every endpoint convention
leave the exact coefficient-one conclusion unchanged.
\end{proof}
\subsection{Random-Intercept Conditional Kernels}
\label{app:random-intercepts}

\begin{lemma}[Measurable conditional kernel for an affine uniform intercept]
\label{lem:conditional-uniform-kernel}
Let $R>0$, let $\beta$ have an arbitrary Borel probability law $\pi$ on
the higher-coordinate cube, with the one-point convention when $d=1$,
and let $U\sim\operatorname{Unif}[-1,1]$ be independent of $\beta$.
Let $G,\rho$ be measurable real-valued functions of $\beta$ and put
$Y=G(\beta)+\rho(\beta)U$.  Suppose $H$ is measurable,
$\pi(H)=1$, and $\rho(\beta)>0$ for every $\beta\in H$.  Then $Y$ has a
jointly measurable regular conditional density $f(a\mid\beta)$ given
$\beta$ such that, for every $\beta\in H$,
\[
f(a\mid\beta)
=\frac{\mathbf1\{|a-G(\beta)|\leq\rho(\beta)\}}
       {2\rho(\beta)},
\qquad
\operatorname*{ess\,sup}_{a\in\mathbb R}f(a\mid\beta)
=\frac{1}{2\rho(\beta)}.
\]
If also $|G(\beta)|+\rho(\beta)\leq R$ on $H$, the density version can
be chosen to vanish outside $[-R,R]$ on every fiber.
\end{lemma}

\begin{proof}
Define the fallback density
\[
g_R(a):=\frac{\mathbf1\{|a|\leq R\}}{2R}
\]
and the piecewise function
\[
f(a\mid\beta)=
\begin{cases}
\displaystyle
\frac{\mathbf1\{|a-G(\beta)|\leq\rho(\beta)\}}
     {2\rho(\beta)},&\beta\in H,\\[8pt]
g_R(a),&\beta\notin H.
\end{cases}
\]
There is no quotient on $H^c$, so this definition contains no hidden
$0/0$ convention.  The set
\[
\{(a,\beta):|a-G(\beta)|\leq\rho(\beta)\}
\]
is measurable, as are $H,G$, and $\rho$.  Strict positivity on $H$
makes $f$ jointly measurable.  Its integral in $a$ equals one on $H$
because the indicated interval has length $2\rho(\beta)$, and equals one
on $H^c$ by the definition of $g_R$.  Thus
\[
Q(A\mid\beta):=\int_A f(a\mid\beta)\,da
\]
is a measurable probability kernel.

Fix Borel sets $A\subseteq\mathbb R$ and $B_0$ in the
higher-coordinate space.  Independence gives the product-law identity
\[
\Pr(Y\in A,\beta\in B_0)
=\int_{B_0}\frac12\int_{-1}^{1}
\mathbf1\{G(b)+\rho(b)u\in A\}\,du\,\pi(db).
\]
For $b\in H$, the positive-width substitution $a=G(b)+\rho(b)u$ gives,
for every nonnegative Borel function $h$,
\[
\frac12\int_{-1}^{1}h(G(b)+\rho(b)u)\,du
=\int_{\mathbb R}h(a)
\frac{\mathbf1\{|a-G(b)|\leq\rho(b)\}}{2\rho(b)}\,da.
\]
Consequently, the inner integral over $A$ equals
$\int_Af(a\mid b)\,da$ on $H$.  Since $H^c$ has $\pi$-measure zero,
replacing the inner integral there by the fallback kernel changes
nothing, and
\[
\Pr(Y\in A,\beta\in B_0)
=\int_{B_0}\int_Af(a\mid b)\,da\,\pi(db).
\]
This rectangle identity proves that $Q$ is a regular conditional law of
the actual $Y$ given the actual $\beta$.  It integrates against $\pi$
itself and needs no Lebesgue density of $\pi$.

For $b\in H$, the displayed density equals the positive constant
$1/[2\rho(b)]$ on an interval of positive length $2\rho(b)$ and is zero
off that interval.  Its essential supremum is therefore exactly
$1/[2\rho(b)]$.  A globally measurable cap version is
\[
K(b)=
\begin{cases}
1/[2\rho(b)],&b\in H,\\
1/(2R),&b\notin H.
\end{cases}
\]
Finally, if $|G|+\rho\leq R$ on $H$, then
$[G(b)-\rho(b),G(b)+\rho(b)]\subseteq[-R,R]$ there, while $g_R$ is also
supported on that interval.  Hence the selected version vanishes outside
$[-R,R]$ on every fiber.
\end{proof}

\begin{proposition}[Reciprocal-width envelope for random intercepts]
\label{prop:random-intercept-envelope}
Under Assumption~\ref{assump:random-intercept-witness} and
Lemma~\ref{lem:conditional-uniform-kernel}, every
$\mu\in\mathcal D_{\mathrm{RI}}$ is supported on $[-R,R]^d$ and has a
jointly measurable regular conditional density, extended by zero outside
$[-R,R]$, satisfying for $\pi_\mu$-almost every $\beta$
\[
f_\mu(a_0\mid\beta)
=\frac{\mathbf1\{|a_0-G_\mu(\beta)|\leq\rho_\mu(\beta)\}}
       {2\rho_\mu(\beta)},
\qquad
K_\mu(\beta)=\frac{1}{2\rho_\mu(\beta)}.
\]
Moreover, $K_\mu$ has a measurable version and
\[
\int K_\mu(\beta)\,\pi_\mu(d\beta)
=\frac12\int\rho_\mu(\beta)^{-1}\,\pi_\mu(d\beta),
\qquad
\bar\kappa_{\mathcal D_{\mathrm{RI}}}
=\frac12M_{\mathrm{RI}}<\infty.
\]
Thus $\mathcal D_{\mathrm{RI}}$ satisfies every component of
Assumption~\ref{assump:averaged-intercept-density}.
\end{proposition}

\begin{proof}
Fix $\mu\in\mathcal D_{\mathrm{RI}}$ and define
\[
H_\mu
:=\{\beta:\rho_\mu(\beta)>0,
             \ |G_\mu(\beta)|+\rho_\mu(\beta)\leq R\}.
\]
This set is measurable, and
Assumption~\ref{assump:random-intercept-witness} gives
$\pi_\mu(H_\mu)=1$.  Apply
Lemma~\ref{lem:conditional-uniform-kernel} with $H=H_\mu$.  It produces
a jointly measurable conditional density with the asserted formula and
cap.  Its support clause gives a version vanishing outside $[-R,R]$ on
every fiber, and its piecewise cap gives a measurable essential-supremum
version on every fiber.

The higher-coordinate vector already lies in $[-R,R]^{d-1}$.  On the
probability-one event $\{\beta\in H_\mu,|U_\mu|\leq1\}$,
\[
|\alpha_0|
\leq|G_\mu(\beta)|+\rho_\mu(\beta)|U_\mu|
\leq|G_\mu(\beta)|+\rho_\mu(\beta)
\leq R.
\]
Therefore the actual coefficient law is supported on $[-R,R]^d$.

The complement of $H_\mu$ is null, so neither the fallback cap nor any
arbitrary value assigned to $\rho_\mu^{-1}$ there contributes.  Hence
\[
\begin{aligned}
\int K_\mu(\beta)\,\pi_\mu(d\beta)
&=\frac12\int_{H_\mu}
  \rho_\mu(\beta)^{-1}\,\pi_\mu(d\beta)\\
&=\frac12\int
  \rho_\mu(\beta)^{-1}\,\pi_\mu(d\beta).
\end{aligned}
\]
Only after this per-law integration do we take the class supremum:
\[
\begin{aligned}
\bar\kappa_{\mathcal D_{\mathrm{RI}}}
&=\sup_{\mu\in\mathcal D_{\mathrm{RI}}}
  \int K_\mu\,d\pi_\mu\\
&=\frac12\sup_{\mu\in\mathcal D_{\mathrm{RI}}}
  \int\rho_\mu^{-1}\,d\pi_\mu
=\frac12M_{\mathrm{RI}}.
\end{aligned}
\]
No essential supremum over $\beta$ has been inserted.  When $d=1$, the
higher-coordinate space is the one-point empty-tuple space and all these
integrals are evaluations at that point.  The kernel, support estimate,
and reciprocal-width identity remain literal.  Joint measurability of
$f_\mu$, measurability of $K_\mu$, and finiteness of the last display
verify all components of
Assumption~\ref{assump:averaged-intercept-density}.
\end{proof}

\begin{proposition}[Uniform root hitting for random intercepts]
\label{prop:random-intercept-root-hitting}
Under Assumption~\ref{assump:random-intercept-witness},
Proposition~\ref{prop:random-intercept-envelope}, and
Proposition~\ref{prop:averaged-root-hitting}, every
$\mu\in\mathcal D_{\mathrm{RI}}$ and
$I\in\mathscr I(\Theta)$ satisfy
\[
\begin{aligned}
\Pr_\mu(Z_\alpha\cap I\neq\varnothing)
&\leq\frac{L_{d,R,\Theta}|I|}{2}
  \int\rho_\mu(\beta)^{-1}\,\pi_\mu(d\beta)\\
&\leq L_{d,R,\Theta}|I|
  \bar\kappa_{\mathcal D_{\mathrm{RI}}}\\
&\leq\frac{L_{d,R,\Theta}M_{\mathrm{RI}}}{2}|I|,
\end{aligned}
\]
and
\[
C_{\mathcal D_{\mathrm{RI}}}
\leq L_{d,R,\Theta}\bar\kappa_{\mathcal D_{\mathrm{RI}}}
\leq\frac12L_{d,R,\Theta}M_{\mathrm{RI}}<\infty.
\]
\end{proposition}

\begin{proof}
Proposition~\ref{prop:random-intercept-envelope} proves from the primitive
random-intercept assumption that the class is cube-supported and satisfies
the full conditional-density and averaged-envelope hypothesis of
Proposition~\ref{prop:averaged-root-hitting}.  Applying that proposition
to $\mathcal D_{\mathrm{RI}}$ gives
\[
\Pr_\mu(Z_\alpha\cap I\neq\varnothing)
\leq L_{d,R,\Theta}|I|\int K_\mu\,d\pi_\mu.
\]
Substitution of the exact per-law identity from
Proposition~\ref{prop:random-intercept-envelope}, followed by the
definition of $M_{\mathrm{RI}}$, gives
\[
\begin{aligned}
L_{d,R,\Theta}|I|\int K_\mu\,d\pi_\mu
&=\frac{L_{d,R,\Theta}|I|}{2}
  \int\rho_\mu^{-1}\,d\pi_\mu\\
&\leq L_{d,R,\Theta}|I|
      \bar\kappa_{\mathcal D_{\mathrm{RI}}}\\
&\leq\frac{L_{d,R,\Theta}M_{\mathrm{RI}}}{2}|I|.
\end{aligned}
\]
The class conclusion follows in the same way.  The fixed law's reciprocal
width is integrated before the class supremum is used.  Since
Proposition~\ref{prop:averaged-root-hitting} is already uniform over all
positive-length intervals, no interval-dependent exceptional set or
probability conversion arises.  The one-point $d=1$ marginal is covered
by the same substitution.
\end{proof}

\begin{proposition}[Exact fixed-width certificate]
\label{prop:fixed-width}
Under Assumption~\ref{assump:random-intercept-witness},
Proposition~\ref{prop:random-intercept-envelope}, and
Proposition~\ref{prop:random-intercept-root-hitting}, suppose
$\mu\in\mathcal D_{\mathrm{RI}}$ has
$\rho_\mu(\beta)=r_\mu$ for $\pi_\mu$-almost every $\beta$, where
$r_\mu>0$.  Then
\[
K_\mu(\beta)=\frac{1}{2r_\mu}
\quad\text{for }\pi_\mu\text{-almost every }\beta,
\qquad
\int K_\mu\,d\pi_\mu=\frac{1}{2r_\mu},
\]
and, for every $I\in\mathscr I(\Theta)$,
\[
\Pr_\mu(Z_\alpha\cap I\neq\varnothing)
\leq\frac{L_{d,R,\Theta}}{2r_\mu}|I|.
\]
If $\mathcal F$ is a nonempty class of such fixed-width members satisfying
the same support condition and
$\sup_{\mu\in\mathcal F}r_\mu^{-1}<\infty$, then
\[
\bar\kappa_{\mathcal F}
=\frac12\sup_{\mu\in\mathcal F}r_\mu^{-1},
\qquad
C_{\mathcal F}
\leq\frac{L_{d,R,\Theta}}2
       \sup_{\mu\in\mathcal F}r_\mu^{-1}<\infty.
\]
\end{proposition}

\begin{proof}
On the full-measure set where the fixed-width identity and
Assumption~\ref{assump:random-intercept-witness} hold,
Proposition~\ref{prop:random-intercept-envelope} gives
\[
K_\mu(\beta)=\frac{1}{2\rho_\mu(\beta)}
=\frac{1}{2r_\mu}.
\]
Integration with respect to the probability law $\pi_\mu$ proves the
cap integral, and substitution into
Proposition~\ref{prop:random-intercept-root-hitting} proves the per-law
root bound.  For a class $\mathcal F$, take the supremum only after the
per-law calculation:
\[
\bar\kappa_{\mathcal F}
=\sup_{\mu\in\mathcal F}\frac{1}{2r_\mu}
=\frac12\sup_{\mu\in\mathcal F}r_\mu^{-1}.
\]
The stated reciprocal-width condition is exactly the class-uniform
integrability condition specialized to fixed widths.  Applying
Proposition~\ref{prop:averaged-root-hitting} to $\mathcal F$ gives the
class conclusion.  For $d=1$, $r_\mu$ is the width of the unique
empty-tuple fiber and all equations remain unchanged.  Together with
Propositions~\ref{prop:random-intercept-envelope}
and~\ref{prop:random-intercept-root-hitting}, this proves the conditional
density, exact cap, averaged budget, root bound, and fixed-width clauses
without any higher-marginal density assumption.
\end{proof}

\subsection{Exact Heteroscedastic Sheet and Conditional Law}
\label{app:heteroscedastic-sheet}

\begin{proposition}[Exact support of the heteroscedastic sheet]
\label{prop:exact-sheet-support}
Fix $d\geq3$, $R>0$, and $q\in(0,1)$, and let $\mu_q$ be the explicit
witness law from the setup.  Define
\[
T_q=[-R/2,R/2]\times[-1,1]
\]
and
\[
F_q:T_q\longrightarrow\mathbb R^d,
\qquad
F_q(z,u)=(\rho_q(z)u,z,z^2/R,0,\ldots,0).
\]
When $d=3$, the trailing list of zeros is empty.  Then
\[
\operatorname{supp}(\mu_q)=F_q(T_q)=S_q,
\]
where
\[
S_q=\left\{(a_0,z,z^2/R,0,\ldots,0):
|z|\leq R/2,\ |a_0|\leq\rho_q(z)\right\}.
\]
The identity includes the collapsed fiber at $z=0$ and the fibers at
$z=\pm R/2$.
\end{proposition}

\begin{proof}
Since $q>0$, the map $z\mapsto|z|^q$, and hence $\rho_q$, is continuous
at every $z$, including zero.  Since $R>0$, all other coordinates of
$F_q$ are continuous.  Thus $F_q$ is continuous on the compact rectangle
$T_q$ and $F_q(T_q)$ is compact, hence closed.

For $(z,u)\in T_q$, $|\rho_q(z)u|\leq\rho_q(z)$, so
$F_q(T_q)\subseteq S_q$.  Conversely, take a point of $S_q$.  If
$z\neq0$, then $\rho_q(z)>0$ and
\[
u=\frac{a_0}{\rho_q(z)}\in[-1,1]
\]
maps to that point.  If $z=0$, then $\rho_q(0)=0$ and the sheet
inequality forces $a_0=0$; choosing $u=0$ again gives the point.  Thus
\[
F_q(T_q)=S_q.
\]

Independence and the uniform laws of $Z$ and $U$ make their joint law
normalized two-dimensional Lebesgue measure on $T_q$.  Every nonempty
relatively open subset of $T_q$, including one meeting an edge or corner,
has positive probability, so the support of $(Z,U)$ is exactly $T_q$.
Moreover, $\alpha^{(q)}=F_q(Z,U)$, and therefore
$\mu_q=(F_q)_\#\mathcal L(Z,U)$.

The pushforward is concentrated on the closed set $F_q(T_q)$, giving one
support inclusion.  For the reverse inclusion, fix $y=F_q(z,u)$ and an
open neighborhood $O\ni y$.  Continuity makes $F_q^{-1}(O)$ a relatively
open subset of $T_q$ containing $(z,u)$, so it has positive joint
probability and
\[
\mu_q(O)=\Pr(F_q(Z,U)\in O)>0.
\]
Thus every image point belongs to the topological support.

At $z=0$, every $(0,u)$ maps to $(0,0,0,\ldots,0)$.  Although
$\Pr(Z=0)=0$, every neighborhood of this point has positive pushforward
mass by the preimage argument, so the collapsed point remains in support.
At $z=\pm R/2$, relative neighborhoods of the latent boundary also have
positive mass; hence the full endpoint fibers, including their
$u=\pm1$ endpoints, remain in support.
\end{proof}

\begin{proposition}[Cube support and almost-sure positive width]
\label{prop:cube-positive-width}
Fix $d\geq3$, $R>0$, and $q\in(0,1)$.  Under
Proposition~\ref{prop:exact-sheet-support},
\[
S_q\subseteq[-R,R]^d,
\qquad
\mu_q([-R,R]^d)=1,
\qquad
\Pr(\rho_q(Z)>0)=1.
\]
These conclusions include $d=3$, the collapsed support point at $z=0$,
and both endpoint fibers $z=\pm R/2$.
\end{proposition}

\begin{proof}
For $|z|\leq R/2$, put $t=2|z|/R$.  Since $R>0$, $t\in[0,1]$, and
$q>0$ gives $t^q\in[0,1]$.  Hence
\[
0\leq\rho_q(z)=\frac R2t^q\leq\frac R2.
\]
Every point of $S_q$ satisfies
\[
|a_0|\leq\rho_q(z)\leq R/2\leq R,
\qquad
|\alpha_1|=|z|\leq R/2\leq R,
\]
and
\[
|\alpha_2|=\frac{z^2}{R}
\leq\frac{(R/2)^2}{R}=\frac R4\leq R.
\]
All remaining coordinates vanish.  This proves
$S_q\subseteq[-R,R]^d$, and
Proposition~\ref{prop:exact-sheet-support} gives cube support of the
actual law.

Because $R>0$ and $q>0$,
\[
\rho_q(z)=0\quad\Longleftrightarrow\quad z=0.
\]
The uniform law on the nondegenerate interval $[-R/2,R/2]$ has no atoms,
so $\Pr(Z=0)=0$ and the almost-sure assertion follows.  This does not
remove the origin from topological support.

When $d=3$, the only coordinates are $a_0,z,z^2/R$, and the same bounds
apply with no trailing coordinates.  At $z=\pm R/2$,
$\rho_q(z)=R/2$ and $z^2/R=R/4$; at $z=0$, all displayed coordinates
vanish.  The upper restriction $q<1$ is not needed for these support
claims; it is used in the cap integral below.
\end{proof}

\begin{lemma}[Recovery of the latent coordinate and inherited independence]
\label{lem:latent-recovery}
Fix $d\geq3$, $R>0$, and $q\in(0,1)$.  For a higher-coordinate vector
$b\in[-R,R]^{d-1}$, let $b_1$ denote its first coordinate and define
\[
G_q(b)=0,
\qquad
\widetilde\rho_q(b)
=\frac R2\left(\frac{2|b_1|}{R}\right)^q,
\qquad
H_q=\{b:0<|b_1|\leq R/2\}.
\]
Then $G_q,\widetilde\rho_q$, and $H_q$ are measurable,
\[
Z=\operatorname{pr}_1(\beta_q(Z)),
\qquad
\sigma(\beta_q(Z))=\sigma(Z),
\]
$U$ is independent of $\beta_q(Z)$,
$\pi_{\mu_q}(H_q)=1$, $\widetilde\rho_q>0$ on $H_q$, and
\[
\rho_q(Z)U
=G_q(\beta_q(Z))+\widetilde\rho_q(\beta_q(Z))U.
\]
Thus all conditions of Lemma~\ref{lem:conditional-uniform-kernel} hold
for this construction.
\end{lemma}

\begin{proof}
The condition $d\geq3$ ensures that the higher-coordinate vector has a
first coordinate.  By definition,
\[
\beta_q(Z)=(Z,Z^2/R,0,\ldots,0),
\]
so projection onto that coordinate recovers $Z$ exactly.  Since
$\beta_q(Z)$ is a measurable function of $Z$,
$\sigma(\beta_q(Z))\subseteq\sigma(Z)$; recovery gives the reverse
inclusion.  The signed linear coordinate prevents any ambiguity from the
quadratic coordinate.

For Borel sets $A\subseteq[-1,1]$ and
$C\subseteq[-R,R]^{d-1}$, independence of $U$ and $Z$ gives
\[
\begin{aligned}
\Pr(U\in A,\beta_q(Z)\in C)
&=\Pr(U\in A,Z\in\beta_q^{-1}(C))\\
&=\Pr(U\in A)\Pr(Z\in\beta_q^{-1}(C))\\
&=\Pr(U\in A)\Pr(\beta_q(Z)\in C).
\end{aligned}
\]
Thus $U$ is independent of the actual conditioning vector.

The functions above are continuous and $H_q$ is Borel.  Along the
actual higher-coordinate curve,
\[
\widetilde\rho_q(\beta_q(z))
=\frac R2\left(\frac{2|z|}{R}\right)^q
=\rho_q(z).
\]
This proves the representation.  Also,
\[
\pi_{\mu_q}(H_q)=\Pr(0<|Z|\leq R/2)=1,
\]
because $Z$ has no atom at zero.  For $b\in H_q$, positivity of $R,q$ and
$b_1\neq0$ imply $\widetilde\rho_q(b)>0$, and
$|G_q(b)|+\widetilde\rho_q(b)\leq R/2\leq R$.  Thus measurability,
representation, independence, full measure, positive width, and support
are all established before Lemma~\ref{lem:conditional-uniform-kernel} is
used.  No reciprocal-width integrability is asserted here.
\end{proof}

\begin{proposition}[Conditional density and exact heteroscedastic cap]
\label{prop:witness-conditional-cap}
Fix $d\geq3$, $R>0$, and $q\in(0,1)$.  Under
Lemma~\ref{lem:latent-recovery} and
Lemma~\ref{lem:conditional-uniform-kernel}, the actual intercept
$\alpha_0^{(q)}=\rho_q(Z)U$ has a jointly measurable regular conditional
density given $\beta_q(Z)$.  One valid version satisfies, for every
$0<|z|\leq R/2$,
\[
f_{\mu_q}(a_0\mid\beta_q(z))
=\frac{\mathbf1\{|a_0|\leq\rho_q(z)\}}{2\rho_q(z)},
\]
and its exact Lebesgue essential-supremum cap is
\[
K_{\mu_q}(\beta_q(z))
=\frac{1}{2\rho_q(z)}
=\frac{1}{R(2|z|/R)^q}.
\]
At $z=0$, neither formula is asserted; the selected conditional version
uses a genuine fallback density on that $\pi_{\mu_q}$-null fiber.
\end{proposition}

\begin{proof}
Lemma~\ref{lem:latent-recovery} proves that
Lemma~\ref{lem:conditional-uniform-kernel} applies with
\[
\beta=\beta_q(Z),\quad G=G_q,\quad
\rho=\widetilde\rho_q,\quad H=H_q,
\quad Y=\alpha_0^{(q)}.
\]
The noise is therefore independent of the actual conditioning vector,
not merely of an unrecovered latent variable.  A valid density version is
\[
f_{\mu_q}(a_0\mid b)=
\begin{cases}
\displaystyle
\frac{\mathbf1\{|a_0|\leq\widetilde\rho_q(b)\}}
     {2\widetilde\rho_q(b)},&b\in H_q,\\[8pt]
\displaystyle
\frac{\mathbf1\{|a_0|\leq R\}}{2R},&b\notin H_q.
\end{cases}
\]
The second line is a probability density because $R>0$.  It defines a
measurable version on the null complement without any quotient involving
$\widetilde\rho_q=0$.

For $0<|z|\leq R/2$, $\beta_q(z)\in H_q$ and
Lemma~\ref{lem:latent-recovery} turns the first line into the asserted
density.  It is the positive constant $1/[2\rho_q(z)]$ on an interval of
positive length $2\rho_q(z)$ and zero elsewhere, so its essential
supremum is exactly that constant.  Since
\[
2\rho_q(z)=R\left(\frac{2|z|}{R}\right)^q,
\]
the cap has the displayed form.

At $z=\pm R/2$, the denominator is $R$ and the cap is exactly $1/R$.
At $z=0$, Proposition~\ref{prop:exact-sheet-support} retains the
collapsed point in topological support, but
$\pi_{\mu_q}(\{\beta_q(0)\})=\Pr(Z=0)=0$.  A regular conditional law is
not determined there, and the fallback is a legal version choice rather
than a positive-mass conditional-uniform assertion.  Combining this
conditional law with Propositions~\ref{prop:exact-sheet-support}
and~\ref{prop:cube-positive-width} gives the exact support, cube,
positive-width, latent-recovery, and nonzero-fiber cap conclusions for the
witness, without evaluating a reciprocal-width integral at the null
fiber.
\end{proof}

\subsection{Unbounded Slice Caps with a Finite Average}
\label{app:witness-cap-integral}

\begin{lemma}[Positive-mass blowup of the heteroscedastic slice cap]
\label{lem:positive-mass-cap}
Fix $d\geq3$, $R>0$, and $q\in(0,1)$.  Under
Proposition~\ref{prop:witness-conditional-cap}, for every $M>0$,
\[
\pi_{\mu_q}\{\beta:K_{\mu_q}(\beta)>M\}>0.
\]
Consequently,
\[
\operatorname*{ess\,sup}_{\beta\sim\pi_{\mu_q}}
K_{\mu_q}(\beta)=\infty.
\]
Both conclusions are unchanged by any legal conditional-density fallback
on the marginal-null fiber $\beta_q(0)$.
\end{lemma}

\begin{proof}
For $z\in[-R/2,R/2]$, define the dimensionless coordinate
\[
t(z):=\frac{2|z|}{R}\in[0,1].
\]
Fix $M>0$ and put
\[
\delta_M:=\min\left\{\frac12,(2RM)^{-1/q}\right\}.
\]
Because $R,q,M$ are positive,
$0<\delta_M\leq1/2$, and
\[
\delta_M^q\leq\frac{1}{2RM},
\qquad
\frac{1}{R\delta_M^q}\geq2M.
\]

Let $b_1$ be the first coordinate of a higher-coordinate vector $b$ and
define the Borel set
\[
A_M:=\left\{b:0<\frac{2|b_1|}{R}<\delta_M\right\}.
\]
Lemma~\ref{lem:latent-recovery} gives $(\beta_q(z))_1=z$.  Since $Z$ is
uniform on an interval of length $R$,
\[
\begin{aligned}
\pi_{\mu_q}(A_M)
&=\Pr\left(0<\frac{2|Z|}{R}<\delta_M\right)\\
&=\frac1R\left|
\left(-\frac{R\delta_M}{2},\frac{R\delta_M}{2}\right)
\setminus\{0\}\right|
=\delta_M>0.
\end{aligned}
\]
For every actual fiber $b=\beta_q(z)\in A_M$, $z\neq0$, and
Proposition~\ref{prop:witness-conditional-cap} gives
\[
K_{\mu_q}(b)
=\frac{1}{Rt(z)^q}
>\frac{1}{R\delta_M^q}
\geq2M>M.
\]
Thus the cap exceeds $M$ on a set of positive marginal measure.

A nonnegative measurable function has finite essential supremum only if
some finite $M$ bounds it outside a null set, which the preceding result
rules out.  The set $A_M$ expressly excludes $b_1=0$, while
$\pi_{\mu_q}\{\beta_q(0)\}=0$, so a fallback modification cannot affect
the conclusion.  At $z=\pm R/2$, the cap is the finite value $1/R$; the
blowup comes solely from positive-mass punctured neighborhoods of zero.
The strict condition $q>0$ is exactly what makes the cap diverge there.
\end{proof}

\begin{lemma}[Exact averaged cap of the power-width witness]
\label{lem:exact-witness-average}
Fix $d\geq3$, $R>0$, and $q\in(0,1)$.  Under
Proposition~\ref{prop:witness-conditional-cap},
\[
\int K_{\mu_q}(\beta)\,\pi_{\mu_q}(d\beta)
=\frac{1}{R(1-q)}<\infty.
\]
For every dimensionless cutoff $\delta\in(0,1]$, the cap mass on
$t(z)\geq\delta$ is
\[
\frac{1-\delta^{1-q}}{R(1-q)},
\]
and the omitted central mass is exactly
$\delta^{1-q}/[R(1-q)]$, which vanishes as $\delta\downarrow0$.
\end{lemma}

\begin{proof}
Retain $t(z)=2|z|/R$ and define
\[
\widehat K_q(z):=
\begin{cases}
\displaystyle\frac{1}{Rt(z)^q},&z\neq0,\\[6pt]
0,&z=0.
\end{cases}
\]
Proposition~\ref{prop:witness-conditional-cap} gives
$\widehat K_q(z)=K_{\mu_q}(\beta_q(z))$ for every $z\neq0$.  They may
differ at zero, but the uniform law assigns zero probability there.
Because $\pi_{\mu_q}=(\beta_q)_\#\mathcal L(Z)$,
\[
\int K_{\mu_q}\,d\pi_{\mu_q}
=\mathbb E[\widehat K_q(Z)].
\]
This remains true even if a selected null-fiber cap is infinite: a
nonnegative extended-valued function supported on a null set has integral
zero.

For integers $n\geq1$, define
\[
g_n(z):=\widehat K_q(z)\mathbf1\{t(z)\geq1/n\}.
\]
Each $g_n$ is nonnegative and Borel.  Its cutoff set increases with $n$,
so $g_n(z)\uparrow\widehat K_q(z)$ for every $z$, including zero.  By
monotone convergence and the exact uniform density $1/R$,
\[
\mathbb E[\widehat K_q(Z)]
=\lim_{n\to\infty}\frac1R
\int_{-R/2}^{R/2}g_n(z)\,dz.
\]

For each $n$, symmetry and the substitution
$t=2z/R$, $dz=(R/2)dt$, yield
\[
\begin{aligned}
\frac1R\int_{-R/2}^{R/2}g_n(z)\,dz
&=\frac{2}{R}\int_{R/(2n)}^{R/2}
  \frac{1}{R(2z/R)^q}\,dz\\
&=\frac1R\int_{1/n}^{1}t^{-q}\,dt\\
&=\frac{1-n^{-(1-q)}}{R(1-q)}.
\end{aligned}
\]
Every normalization is visible: the first $1/R$ is the density of $Z$,
the second is part of the cap, symmetry contributes two, and
$dz=(R/2)dt$ leaves the single prefactor $1/R$.

Since $q<1$, $1-q>0$ and $n^{-(1-q)}\to0$, proving the exact integral.
Replacing $1/n$ by any $\delta\in(0,1]$ gives the cutoff mass; subtracting
it from the full integral gives the omitted mass.

The point $z=0$ is excluded only from the pointwise formula and remains a
null point in the integration domain.  At $z=\pm R/2$, $t=1$ and the cap
is $1/R$, so no endpoint singularity occurs.  At the excluded lower
boundary $q=0$, the mean remains finite but
Lemma~\ref{lem:positive-mass-cap} fails; at the excluded upper boundary
$q=1$, the logarithmic integral $\int_0^1t^{-1}dt$ diverges.  The degree
condition $d\geq3$ identifies the coefficient vector but introduces no
normalization factor.
\end{proof}

\begin{proposition}[Root hitting for the heteroscedastic witness]
\label{prop:witness-root-hitting}
Fix $d\geq3$, $R>0$, $q\in(0,1)$, and compact
$\Theta\subseteq\mathbb R$.  Under
Proposition~\ref{prop:cube-positive-width},
Proposition~\ref{prop:witness-conditional-cap},
Lemma~\ref{lem:exact-witness-average}, and
Proposition~\ref{prop:averaged-root-hitting}, the singleton class
$\{\mu_q\}$ satisfies
Assumption~\ref{assump:averaged-intercept-density}, with
\[
\bar\kappa_{\{\mu_q\}}
=\int K_{\mu_q}(\beta)\,\pi_{\mu_q}(d\beta)
=\frac{1}{R(1-q)}.
\]
Consequently, for every $I\in\mathscr I(\Theta)$,
\[
\Pr_{\mu_q}(Z_\alpha\cap I\neq\varnothing)
\leq\frac{L_{d,R,\Theta}|I|}{R(1-q)},
\qquad
C_{\{\mu_q\}}
\leq\frac{L_{d,R,\Theta}}{R(1-q)}.
\]
\end{proposition}

\begin{proof}
Proposition~\ref{prop:cube-positive-width} proves cube support.
Proposition~\ref{prop:witness-conditional-cap} supplies a jointly
measurable regular conditional-density version of the actual intercept
given the actual higher-coordinate vector, together with a measurable
essential-supremum cap.  On the full-measure set of nonzero fibers, that
cap has the formula used in Lemma~\ref{lem:exact-witness-average}, which
proves the cap integral is finite and has the displayed exact value.
A kernel or cap modification at $\beta_q(0)$ changes neither the integral
nor any conditional identity on the full-measure set.

Thus the singleton class satisfies every component of
Assumption~\ref{assump:averaged-intercept-density}; none is assumed for
the witness.  Conditional-density and cap measurability come from
Proposition~\ref{prop:witness-conditional-cap}, and finite averaged
envelope comes from Lemma~\ref{lem:exact-witness-average}.  A supremum
over a singleton is its sole value.

The conditions of Proposition~\ref{prop:averaged-root-hitting} are now
verified.  Applying its per-law conclusion and substituting the exact
cap integral gives the interval inequality.  Division by $|I|>0$ and
the defining interval supremum give the bound on $C_{\{\mu_q\}}$.  There
is no class-supremum loss and no pointwise-in-$\beta$ cap assumption.
Finiteness follows from $R>0$, $1-q>0$, and
Lemma~\ref{lem:derivative-envelope}.  Combining this proposition with
Lemma~\ref{lem:positive-mass-cap} proves simultaneously that the slice
cap is essentially unbounded, its mean is exactly finite, and the
root-hitting constant has the stated exact dependence, with no cutoff
remainder or hidden factor.
\end{proof}

\subsection{Geometry Beyond the Full-Rank Affine-Latent Baseline}
\label{app:affine-exclusion}

\begin{proposition}[Exact affine hull of the witness sheet]
\label{prop:sheet-affine-hull}
Fix $d\geq3$, $R>0$, and $q\in(0,1)$.  Under
Proposition~\ref{prop:exact-sheet-support}, define
\[
E_d:=\{(a,z,w,0,\ldots,0):a,z,w\in\mathbb R\}
\subseteq\mathbb R^d.
\]
Then
\[
\operatorname{aff}(S_q)=E_d,
\qquad
\dim\operatorname{aff}(S_q)=3.
\]
For $d=3$, $E_d=\mathbb R^3$.
\end{proposition}

\begin{proof}
Every point of $S_q$ has only its first three coordinates possibly
nonzero, and the collapsed point
\[
o=(0,0,0,0,\ldots,0)
\]
belongs to $S_q$.  Hence $\operatorname{aff}(S_q)$ is a linear subspace
contained in $E_d$.

At either endpoint $z=\pm R/2$,
\[
\rho_q(\pm R/2)
=\frac R2\left(\frac{2(R/2)}R\right)^q
=\frac R2.
\]
Therefore the exact support contains
\[
o,
\qquad
p_+=\left(0,\frac R2,\frac R4,0,\ldots,0\right),
\qquad
p_-=\left(0,-\frac R2,\frac R4,0,\ldots,0\right),
\]
and
\[
r=\left(\frac R2,\frac R2,\frac R4,0,\ldots,0\right).
\]
The projections of $o,p_+,p_-$ onto the
$(\alpha_1,\alpha_2)$ coordinates are
\[
(0,0),\qquad(R/2,R/4),\qquad(-R/2,R/4),
\]
and the determinant of the two nonzero vectors is
\[
\det\begin{pmatrix}R/2&-R/2\\R/4&R/4\end{pmatrix}
=\frac{R^2}{4}>0.
\]
Thus the curved projection supplies two independent affine directions.
More explicitly, the direction space contains
\[
p_+-p_-=Re_1,
\qquad
p_++p_-=\frac R2e_2,
\qquad
r-p_+=\frac R2e_0.
\]
Because $R>0$, these span $E_d$.  This proves the reverse containment and
the affine-hull identity.  For $d=3$, the same directions span all of
$\mathbb R^3$.  The collapsed origin causes no loss because it belongs
to the topological support, while the nonzero directions come from the
endpoint fibers.
\end{proof}

\begin{proposition}[Zero three-volume inside the affine hull]
\label{prop:sheet-null-volume}
Fix $d\geq3$, $R>0$, and $q\in(0,1)$.  Under
Proposition~\ref{prop:exact-sheet-support}, define
\[
D_q:=\{(a,z)\in\mathbb R^2:|z|\leq R/2,
\ |a|\leq\rho_q(z)\}
\]
and
\[
T:D_q\longrightarrow E_d,
\qquad
T(a,z)=(a,z,z^2/R,0,\ldots,0).
\]
Then $T(D_q)=S_q$, the map $T$ is $\sqrt2$-Lipschitz in Euclidean norm,
and, under Proposition~\ref{prop:sheet-affine-hull},
\[
\mathcal H^3_{E_d}(S_q)=0.
\]
\end{proposition}

\begin{proof}
The definitions give both inclusions directly: each $(a,z)\in D_q$ maps
to the displayed sheet, and each sheet point is $T(a,z)$ for its first
two coordinates.  Hence
\[
T(D_q)=S_q.
\]
No division by $\rho_q(z)$ occurs, including at $z=0$.

Take $(a,z),(a',z')\in D_q$.  Since
$|z|,|z'|\leq R/2$,
\[
\left|\frac{z^2-(z')^2}{R}\right|
=\frac{|z-z'||z+z'|}{R}
\leq|z-z'|.
\]
Consequently,
\[
\begin{aligned}
\|T(a,z)-T(a',z')\|_2^2
&=|a-a'|^2+|z-z'|^2
 +\left|\frac{z^2-(z')^2}{R}\right|^2\\
&\leq|a-a'|^2+2|z-z'|^2\\
&\leq2\bigl(|a-a'|^2+|z-z'|^2\bigr).
\end{aligned}
\]
This is the claimed $\sqrt2$-Lipschitz estimate.

Proposition~\ref{prop:cube-positive-width} gives
$\rho_q(z)\leq R/2$, so $D_q$ lies in the square
$[-R/2,R/2]^2$.  A mesh of side length $\delta$ covers this square with
at most $(\lceil R/\delta\rceil+1)^2$ squares, each of diameter at most
$\sqrt2\delta$.  The sum of cubed diameters is bounded by
\[
(\lceil R/\delta\rceil+1)^2(\sqrt2\delta)^3,
\]
which tends to zero as $\delta\downarrow0$.  Thus
$\mathcal H^3(D_q)=0$.  An $L$-Lipschitz map satisfies
$\mathcal H^s(F(A))\leq L^s\mathcal H^s(A)$: map a fine cover of $A$,
noting that each diameter grows by at most $L$, and then take the
defining infimum and limit.  Applying this cover estimate gives
\[
\mathcal H^3_{E_d}(S_q)
=\mathcal H^3_{E_d}(T(D_q))
\leq(\sqrt2)^3\mathcal H^3(D_q)=0.
\]

The regular map here is $T(a,z)$.  The latent parametrization
$(z,u)\mapsto(\rho_q(z)u,z,z^2/R,0,\ldots,0)$ is not used; its first
coordinate is only Holder at zero when $q<1$.  This distinction handles
the collapsed fiber and the minimal case $d=3$ without any unsupported
regularity assertion.
\end{proof}

\begin{lemma}[Fixed-monic descending-coordinate embedding]
\label{lem:monic-embedding}
Under Propositions~\ref{prop:sheet-affine-hull}
and~\ref{prop:sheet-null-volume}, define
\[
\iota_d:\mathbb R^d\longrightarrow\mathbb R^{d+1},
\qquad
\iota_d(a_0,a_1,\ldots,a_{d-1})
=(1,a_{d-1},a_{d-2},\ldots,a_0).
\]
Then $\iota_d$ places the branch coefficient vector in descending full-
coefficient order with fixed leading coefficient $\alpha_d=1$.  It is an
affine isometry, and
\[
\operatorname{supp}((\iota_d)_\#\mu_q)=\iota_d(S_q),
\]
\[
\dim\operatorname{aff}(\iota_d(S_q))=3,
\qquad
\mathcal H^3_{\operatorname{aff}(\iota_d(S_q))}
  (\iota_d(S_q))=0.
\]
\end{lemma}

\begin{proof}
For $x,y\in\mathbb R^d$, the leading coordinates of $\iota_d(x)$ and
$\iota_d(y)$ cancel and the remaining coordinates are merely reversed.
Thus
\[
\|\iota_d(x)-\iota_d(y)\|_2=\|x-y\|_2.
\]
The map is an affine isometry onto the closed fixed-monic hyperplane.  Its
formula gives the exact ordering bridge from
$(\alpha_0,\ldots,\alpha_{d-1})$ to
$(\alpha_d,\ldots,\alpha_0)$ with $\alpha_d=1$.

Write $\iota_d(x)=c_*+Q_dx$.  The isometry gives
$Q_d^\top Q_d=I_d$.  If a lower-coordinate affine map has linear part
$B_0$, then its fixed-monic composition has linear part $Q_dB_0$, and
\[
(Q_dB_0)^\top(Q_dB_0)=B_0^\top B_0.
\]
Thus the translation preserves full column rank and the exact Gram-
determinant normalization.

Because $\iota_d$ is a homeomorphism onto its closed image, it transports
support exactly: every neighborhood of $\iota_d(x)$ with $x\in S_q$
pulls back to a neighborhood of a support point and has positive
pushforward mass, while the pushforward is concentrated on the closed set
$\iota_d(S_q)$.  Affine isometries commute with affine hulls and preserve
their dimensions and induced Hausdorff volumes.  The remaining claims
follow from Propositions~\ref{prop:sheet-affine-hull}
and~\ref{prop:sheet-null-volume}.
\end{proof}

The affine generalization in Appendix Theorem 18 of
\citet{balcan2020semibandit} uses the descending full coefficient vector,
a full-column-rank linear part $A$, affine-image volume proportional to
$\sqrt{\det(A^\top A)}$, and an induced density divided by the same
factor.  In particular, a latent cube $[-R_{\mathrm{lat}},R_{\mathrm{lat}}]^k$
has image volume
$\sqrt{\det(A^\top A)}(2R_{\mathrm{lat}})^k$, and a latent density bounded
by $\kappa$ induces a within-image density bounded by
$\kappa/\sqrt{\det(A^\top A)}$.  The proof in that source unambiguously
uses the full descending vector $(\alpha_d,\ldots,\alpha_0)$; this is the
convention reconciled by Lemma~\ref{lem:monic-embedding}.  The source's
no-forced-root premise belongs to its subsequent root-hitting conclusion
and is not needed for the representability question considered here.  The
following proposition derives the exact measure statement directly.

\begin{proposition}[Full-rank affine pushforward geometry]
\label{prop:affine-pushforward}
Let $k\geq1$.  Let $X\in\mathbb R^k$ have a probability density $p$
with respect to $\mathcal L^k$, with $p\leq\kappa<\infty$ almost
everywhere and bounded support.  Let
\[
\mathfrak f(x)=b+Ax,
\qquad
A\in\mathbb R^{(d+1)\times k},
\qquad
\operatorname{rank}(A)=k,
\]
and put
\[
L=b+\operatorname{im}(A),
\qquad
J_A:=\sqrt{\det(A^\top A)}>0.
\]
Then $\nu:=\mathfrak f_\#\mathcal L(X)$ is absolutely continuous with
respect to $\mathcal H_L^k$, with exact density
\[
g(b+Ax)=\frac{p(x)}{J_A}
\quad\text{for Lebesgue-almost every }x,
\qquad
g\leq\frac{\kappa}{J_A}.
\]
Moreover,
\[
\operatorname{aff}(\operatorname{supp}\nu)=L,
\qquad
\dim\operatorname{aff}(\operatorname{supp}\nu)=k.
\]
\end{proposition}

\begin{proof}
Full column rank makes $\mathfrak f$ one-to-one from $\mathbb R^k$ onto
$L$, with inverse on $L$
\[
x=A^\dagger(y-b),
\qquad
A^\dagger=(A^\top A)^{-1}A^\top.
\]
The $k$-dimensional volume scaling of $A$ is the Gram determinant:
for every Borel set $B_0\subseteq\mathbb R^k$,
\[
\mathcal H_L^k(b+AB_0)=J_A\mathcal L^k(B_0).
\]
Applying this identity first to simple functions and then to the
nonnegative density $p$ gives, for every Borel $C\subseteq L$,
\[
\begin{aligned}
\nu(C)
&=\int_{\mathfrak f^{-1}(C)}p(x)\,dx\\
&=\int_C
\frac{p(A^\dagger(y-b))}{J_A}\,d\mathcal H_L^k(y).
\end{aligned}
\]
This proves the exact density, its bound, and within-image absolute
continuity.

It remains to justify the support dimension because bounded density does
not require $p$ to be positive throughout a bounding cube.  Every Borel
probability on Euclidean space gives full mass to its topological support:
the complement is the union of a countable base of open sets of measure
zero.  If
$H=\operatorname{aff}(\operatorname{supp}\nu)$ were a proper affine
subspace of $L$, then $\mathcal H_L^k(H)=0$ but $\nu(H)=1$, contradicting
the absolute continuity just proved.  Hence $H=L$, with no full-support or
positive lower-density assumption on $X$.
\end{proof}

\begin{proposition}[Exclusion of every affine-latent dimension]
\label{prop:affine-nonmembership}
Fix $d\geq3$, $R>0$, and $q\in(0,1)$.  Under
Proposition~\ref{prop:exact-sheet-support},
Lemma~\ref{lem:monic-embedding}, and
Proposition~\ref{prop:affine-pushforward}, the law
$(\iota_d)_\#\mu_q$ is not the affine pushforward of any bounded-support,
bounded-Lebesgue-density latent probability under a full-column-rank
affine map into $\mathbb R^{d+1}$.  Consequently, $\mu_q$ is not an
affine image of such a latent law in the fixed-monic affine convention of
\citet{balcan2020semibandit}.
\end{proposition}

\begin{proof}
If $\mu_q$ had such a representation in ascending lower-coordinate order,
composition with $\iota_d$ would give one in descending full-coordinate
order.  Lemma~\ref{lem:monic-embedding} shows that this composition
preserves full column rank and the Gram determinant.  It therefore
suffices to exclude the descending-order representations.

Suppose for contradiction that for some positive integer $k$,
\[
(\iota_d)_\#\mu_q=\mathfrak f_\#\mathcal L(X),
\]
where $X$ has a bounded Lebesgue density in $\mathbb R^k$ and
$\mathfrak f(x)=b+Ax$ has full-column-rank $A$.  Equality of laws gives
equality of topological supports.  By Lemma~\ref{lem:monic-embedding},
the left support has affine dimension three.  By
Proposition~\ref{prop:affine-pushforward}, the right support has affine
dimension $k$.  Thus
\[
k=3.
\]
Every positive latent dimension $k\neq3$, including $k=1,2$ and every
$k>3$ compatible with full column rank, is excluded by dimension alone.

For $k=3$, let $L=b+\operatorname{im}(A)$.  Proposition~\ref{prop:affine-pushforward}
gives
\[
\mathfrak f_\#\mathcal L(X)\ll\mathcal H_L^3,
\qquad
L=\operatorname{aff}(\operatorname{supp}
(\mathfrak f_\#\mathcal L(X))).
\]
Equality of laws and Lemma~\ref{lem:monic-embedding} identify $L$ with
$\operatorname{aff}(\iota_d(S_q))$ and give
\[
\mathcal H_L^3(\iota_d(S_q))=0.
\]
Absolute continuity implies
\[
(\mathfrak f_\#\mathcal L(X))(\iota_d(S_q))=0,
\]
whereas exact support and the embedding give
\[
((\iota_d)_\#\mu_q)(\iota_d(S_q))=\mu_q(S_q)=1,
\]
a contradiction.  This is the missing dimension-three argument that
curvature alone would not supply.

If a zero-dimensional latent variable is admitted, its affine image is a
single point and has affine-support dimension zero, again contradicting
Lemma~\ref{lem:monic-embedding}.  If $k>d+1$, no
$(d+1)\times k$ matrix has full column rank.  Thus every nonnegative
latent dimension is covered.  In the minimal case $d=3$, the branch
sheet spans all three lower-coefficient coordinates, its monic embedding
lies in $\mathbb R^4$, and the same $k=3$ null-volume contradiction
applies.

The conclusion does not address rank-deficient affine maps paired with
singular latent laws or any broader comparison class.  It uses neither a
no-forced-root condition nor any root-section estimate or root bound from
the cited source.  Propositions~\ref{prop:sheet-affine-hull}
and~\ref{prop:sheet-null-volume}, together with the exact monic embedding
and affine-pushforward density calculation, provide both the dimension
and within-hull-volume arguments needed for this precisely limited
nonmembership result.
\end{proof}

\subsection{The Sufficient Theorem and Baseline Reduction}
\label{app:theorem-composition}

\begin{proposition}[Complete averaged-intercept sufficient conclusion]
\label{prop:complete-sufficient-theorem}
Fix the basic setting.  The derivative envelope satisfies
\[
L_{d,R,\Theta}
\leq dB^{d-1}+R\sum_{k=1}^{d-1}kB^{k-1}<\infty,
\]
with the sum zero when $d=1$.  Under
Assumption~\ref{assump:averaged-intercept-density}, simultaneously for
every $\mu\in\mathcal D$ and $I\in\mathscr I(\Theta)$,
\[
\Pr_\mu(Z_\alpha\cap I\neq\varnothing)
\leq L_{d,R,\Theta}|I|
      \int K_\mu(\beta)\,\pi_\mu(d\beta)
\leq\bar\kappa_{\mathcal D}L_{d,R,\Theta}|I|,
\]
and
\[
C_{\mathcal D}
\leq\bar\kappa_{\mathcal D}L_{d,R,\Theta}<\infty.
\]

Separately, under Assumption~\ref{assump:random-intercept-witness}, every
$\mu\in\mathcal D_{\mathrm{RI}}$ is cube-supported and, for
$\pi_\mu$-almost every $\beta$,
\[
f_\mu(a_0\mid\beta)
=\frac{\mathbf1\{|a_0-G_\mu(\beta)|\leq\rho_\mu(\beta)\}}
       {2\rho_\mu(\beta)},
\qquad
K_\mu(\beta)=\frac1{2\rho_\mu(\beta)}.
\]
For every such law and every $I\in\mathscr I(\Theta)$,
\[
\bar\kappa_{\mathcal D_{\mathrm{RI}}}=\frac12M_{\mathrm{RI}},
\qquad
\Pr_\mu(Z_\alpha\cap I\neq\varnothing)
\leq\frac{L_{d,R,\Theta}|I|}{2}
      \int\rho_\mu^{-1}\,d\pi_\mu,
\]
and
\[
C_{\mathcal D_{\mathrm{RI}}}
\leq\frac12L_{d,R,\Theta}M_{\mathrm{RI}}<\infty.
\]
If $\rho_\mu=r_\mu>0$ almost surely, then
\[
K_\mu=\frac1{2r_\mu}\quad\pi_\mu\text{-almost surely},
\qquad
\Pr_\mu(Z_\alpha\cap I\neq\varnothing)
\leq\frac{L_{d,R,\Theta}}{2r_\mu}|I|.
\]
For every nonempty class $\mathcal F$ of such fixed-width members obeying
the support condition and
$\sup_{\mu\in\mathcal F}r_\mu^{-1}<\infty$,
\[
\bar\kappa_{\mathcal F}
=\frac12\sup_{\mu\in\mathcal F}r_\mu^{-1},
\qquad
C_{\mathcal F}
\leq\frac{L_{d,R,\Theta}}2
      \sup_{\mu\in\mathcal F}r_\mu^{-1}<\infty.
\]

For every $d\geq3$, $R>0$, and $q\in(0,1)$,
\[
\operatorname{supp}(\mu_q)
=\left\{(a_0,z,z^2/R,0,\ldots,0):
|z|\leq R/2,\ |a_0|\leq\rho_q(z)\right\}
\subseteq[-R,R]^d,
\qquad
\Pr(\rho_q(Z)>0)=1.
\]
For $a_0\in\mathbb R$ and $0<|z|\leq R/2$, one conditional-density
version satisfies
\[
f_{\mu_q}(a_0\mid\beta_q(z))
=\frac{\mathbf1\{|a_0|\leq\rho_q(z)\}}{2\rho_q(z)},
\qquad
K_{\mu_q}(\beta_q(z))
=\frac1{R(2|z|/R)^q}.
\]
Moreover,
\[
\operatorname*{ess\,sup}_{\beta\sim\pi_{\mu_q}}
K_{\mu_q}(\beta)=\infty,
\qquad
\int K_{\mu_q}\,d\pi_{\mu_q}=\frac1{R(1-q)},
\]
and, for every $I\in\mathscr I(\Theta)$,
\[
\Pr_{\mu_q}(Z_\alpha\cap I\neq\varnothing)
\leq\frac{L_{d,R,\Theta}|I|}{R(1-q)},
\qquad
C_{\{\mu_q\}}
\leq\frac{L_{d,R,\Theta}}{R(1-q)}.
\]
The support has affine-hull dimension three and zero induced
three-dimensional Hausdorff volume there, and $\mu_q$ is not a
full-column-rank affine pushforward of a bounded-support,
bounded-Lebesgue-density latent law under the fixed-monic convention.
This last conclusion does not extend by assertion to rank-deficient maps
or singular latent laws.

All displayed conclusions are one-way and sufficient.  They assert no
necessity, converse, indexed uniform polynomial dependence on $(d,R)$,
root simplicity, root separation, or transversality, and supply neither a
learner nor an online guarantee.
\end{proposition}

\begin{proof}
Lemma~\ref{lem:derivative-envelope} gives the first display, including the
empty-sum case.  Proposition~\ref{prop:averaged-root-hitting} gives the
general per-law inequality and closes both defining suprema for
$C_{\mathcal D}$, without any mode conversion.

Under the separately scoped random-intercept assumption,
Proposition~\ref{prop:random-intercept-envelope} proves cube support, the
conditional density, exact cap, per-law cap integral, and exact class
envelope.  It thereby verifies
Assumption~\ref{assump:averaged-intercept-density} for
$\mathcal D_{\mathrm{RI}}$ rather than assuming it.
Proposition~\ref{prop:random-intercept-root-hitting} gives the per-law and
class root bounds, and Proposition~\ref{prop:fixed-width} gives the
fixed-width statements under precisely the class-uniform reciprocal-width
condition.  None of these results needs a density or coordinate
independence assumption for $\pi_\mu$.

For the explicit witness,
Proposition~\ref{prop:exact-sheet-support} gives the support equality,
Proposition~\ref{prop:cube-positive-width} gives cube containment and
almost-sure positivity, and
Proposition~\ref{prop:witness-conditional-cap} gives the selected
conditional density and exact nonzero-fiber cap.  The collapsed zero
fiber remains in topological support but has zero conditioning-marginal
mass and is assigned a genuine fallback density.
Lemma~\ref{lem:positive-mass-cap} gives the infinite essential supremum,
Lemma~\ref{lem:exact-witness-average} gives the exact finite integral, and
Proposition~\ref{prop:witness-root-hitting} gives the interval and
singleton bounds.

Propositions~\ref{prop:sheet-affine-hull}
and~\ref{prop:sheet-null-volume} give the two independent structural
facts.  Proposition~\ref{prop:affine-nonmembership} applies them only to
the full-column-rank latent-Lebesgue-density convention, so the conclusion
does not rely on curvature alone and does not overstate the excluded
class.

The final scope paragraph records the direction of the established
implications.  None of the cited propositions supplies a reverse
implication, a polynomial-family bound, or a root-regularity conclusion.
Thus every displayed rate retains its exact exposed parameters and no
unstated inference is added.
\end{proof}

\begin{lemma}[A bounded joint density implies the averaged intercept cap]
\label{lem:joint-density-bridge}
Fix the basic cube setting and let $\mathcal A$ be a nonempty class of
Borel probability laws on $[-R,R]^d$.  Suppose there is a common
$\kappa_{\mathrm{joint}}\in(0,\infty)$ such that every
$\nu\in\mathcal A$ has a Borel density $h_\nu(a_0,\beta)$ with respect
to $\mathcal L^d$, extended by zero outside the cube, and
\[
0\leq h_\nu(a_0,\beta)
\leq\kappa_{\mathrm{joint}}
\quad\text{for Lebesgue-almost every }(a_0,\beta).
\]
Then $\mathcal A$ satisfies
Assumption~\ref{assump:averaged-intercept-density}, and one may choose
measurable conditional-density and cap versions such that
\[
\int K_\nu(\beta)\,\pi_\nu(d\beta)
\leq\kappa_{\mathrm{joint}}(2R)^{d-1}
\quad\text{for every }\nu\in\mathcal A.
\]
Consequently,
\[
\bar\kappa_{\mathcal A}
\leq\kappa_{\mathrm{joint}}(2R)^{d-1}.
\]
For $d=1$, $(2R)^0=1$ and the same statements hold.
\end{lemma}

\begin{proof}
First suppose $d\geq2$ and fix $\nu\in\mathcal A$.  Define the
higher-coordinate marginal density
\[
g_\nu(\beta):=\int_{\mathbb R}h_\nu(a_0,\beta)\,da_0.
\]
Tonelli's theorem makes $g_\nu$ Borel and gives
$\pi_\nu(d\beta)=g_\nu(\beta)d\beta$.  Because $h_\nu$ is supported on
the cube and bounded by $\kappa_{\mathrm{joint}}$,
\[
0\leq g_\nu(\beta)\leq2R\kappa_{\mathrm{joint}}
\quad\text{for Lebesgue-almost every }\beta,
\qquad
\int g_\nu(\beta)d\beta=1.
\]
Thus the measurable set
\[
G_\nu:=\{\beta:0<g_\nu(\beta)<\infty\}
\]
has full $\pi_\nu$-measure: $\{g_\nu=0\}$ has zero marginal mass, and
$\{g_\nu=\infty\}$ is Lebesgue-null by integrability and also has zero
marginal mass.

Define the sectionwise essential supremum
\[
H_\nu(\beta)
:=\operatorname*{ess\,sup}_{a_0\in\mathbb R}
h_\nu(a_0,\beta).
\]
It has a measurable extended-valued version.  Indeed, for every rational
$t\geq0$, Tonelli makes
\[
m_t(\beta):=\int_{-R}^{R}
\mathbf1\{h_\nu(a_0,\beta)>t\}\,da_0
\]
measurable, and
$\{H_\nu>t\}=\{m_t>0\}$.  The rational superlevel sets determine a
measurable version of $H_\nu$.  Fubini applied to the joint bound gives
\[
H_\nu(\beta)\leq\kappa_{\mathrm{joint}}
\quad\text{for Lebesgue-almost every }\beta.
\]

Now define
\[
f_\nu(a_0\mid\beta)=
\begin{cases}
h_\nu(a_0,\beta)/g_\nu(\beta),&\beta\in G_\nu,\\[4pt]
\mathbf1\{|a_0|\leq R\}/(2R),&\beta\notin G_\nu.
\end{cases}
\]
This is jointly measurable and makes no division on a zero or infinite
marginal-density fiber.  Each line integrates to one: on $G_\nu$ this is
the definition of $g_\nu$, and on $G_\nu^c$ it follows from $R>0$.
For Borel sets $A\subseteq\mathbb R$ and
$D\subseteq\mathbb R^{d-1}$, Tonelli gives
\[
\begin{aligned}
\int_D\int_A f_\nu(a_0\mid\beta)\,da_0\,\pi_\nu(d\beta)
&=\int_{D\cap G_\nu}\int_A
 h_\nu(a_0,\beta)\,da_0\,d\beta\\
&=\nu\{\alpha_0\in A,\ \beta\in D\}.
\end{aligned}
\]
The omitted part $D\cap G_\nu^c$ has zero marginal and joint-law mass.
Thus $f_\nu$ is a jointly measurable regular conditional density of the
actual intercept given the actual higher coordinates, extended by zero
outside $[-R,R]$.

Its measurable essential-supremum version is
\[
K_\nu(\beta)=
\begin{cases}
H_\nu(\beta)/g_\nu(\beta),&\beta\in G_\nu,\\[4pt]
1/(2R),&\beta\notin G_\nu.
\end{cases}
\]
On $G_\nu$, multiplication by the finite positive number
$g_\nu(\beta)$ commutes with the Lebesgue essential supremum, and hence,
for Lebesgue-almost every $\beta\in G_\nu$,
\[
g_\nu(\beta)K_\nu(\beta)
=H_\nu(\beta)
\leq\kappa_{\mathrm{joint}}.
\]
The complement of $G_\nu$ is marginal-null, so its fallback cap
contributes nothing.  Integration over the exact higher-coordinate cube
gives
\[
\begin{aligned}
\int K_\nu(\beta)\,\pi_\nu(d\beta)
&=\int_{G_\nu\cap[-R,R]^{d-1}}
  g_\nu(\beta)K_\nu(\beta)\,d\beta\\
&\leq\kappa_{\mathrm{joint}}
  \operatorname{Leb}^{d-1}([-R,R]^{d-1})\\
&=\kappa_{\mathrm{joint}}(2R)^{d-1}.
\end{aligned}
\]
This establishes joint measurability of the conditional density,
measurability of its cap, and finite cap integral for $d\geq2$.

When $d=1$, the conditioning space is the one-point empty-tuple space.
The conditional density is the one-dimensional joint density
$h_\nu(a_0)$ and
\[
K_\nu
=\operatorname*{ess\,sup}_{a_0}h_\nu(a_0)
\leq\kappa_{\mathrm{joint}}
=\kappa_{\mathrm{joint}}(2R)^0.
\]
The cap integral is evaluation at that one point.  Taking the class
supremum proves the asserted envelope bound.  The density and cap
constructions verify every part of
Assumption~\ref{assump:averaged-intercept-density} for $\mathcal A$;
the joint-density premise is confined to this comparison.
\end{proof}

\begin{proposition}[Recovery of the monic bounded-joint-density baseline]
\label{prop:joint-density-baseline}
Under the basic setting and the hypotheses of
Lemma~\ref{lem:joint-density-bridge}, that lemma and
Proposition~\ref{prop:averaged-root-hitting} imply that every
$\nu\in\mathcal A$ and $I\in\mathscr I(\Theta)$ satisfy
\[
\Pr_\nu(Z_\alpha\cap I\neq\varnothing)
\leq\kappa_{\mathrm{joint}}(2R)^{d-1}
      L_{d,R,\Theta}|I|,
\]
and
\[
C_{\mathcal A}
\leq\kappa_{\mathrm{joint}}(2R)^{d-1}L_{d,R,\Theta}.
\]
If $B=1$, then
\[
C_{\mathcal A}
\leq\kappa_{\mathrm{joint}}(2R)^{d-1}
\left(d+\frac{Rd(d-1)}2\right).
\]
For $d=1$, $(2R)^0=1$ and the bracket equals one.
\end{proposition}

\begin{proof}
Lemma~\ref{lem:joint-density-bridge} constructs the required conditional
density and cap versions for $\mathcal A$, verifies
Assumption~\ref{assump:averaged-intercept-density}, and proves
\[
\bar\kappa_{\mathcal A}
\leq\kappa_{\mathrm{joint}}(2R)^{d-1}.
\]
Proposition~\ref{prop:averaged-root-hitting} therefore applies to the
comparison class.  Its per-law and class conclusions yield the first two
displays without using the joint-density premise anywhere in
Proposition~\ref{prop:complete-sufficient-theorem}.

Lemma~\ref{lem:derivative-envelope} gives
\[
L_{d,R,\Theta}
\leq dB^{d-1}+R\sum_{k=1}^{d-1}kB^{k-1}.
\]
When $B=1$, every power equals one and
\[
\sum_{k=1}^{d-1}k=\frac{d(d-1)}2,
\]
with both sides zero for $d=1$.  Substitution gives the final display.
Every factor is exact: conditioning on the intercept leaves the
higher-coordinate cube volume $(2R)^{d-1}$, and the monic derivative
contributes $d+Rd(d-1)/2$.  Thus the $B=1$ statement recovers the monic
ambient-volume scale rather than an asymptotic surrogate.  In the absence
of the optional joint-density premise, the general averaged-envelope
theorem and both constructive specializations remain unchanged.
\end{proof}

\subsection{Proof of the Main Theorem}

\begin{proof}[Proof of Theorem~\ref{thm:main}]
Lemma~\ref{lem:derivative-envelope} proves the displayed deterministic
bound on $L_{d,R,\Theta}$.  Under
Assumption~\ref{assump:averaged-intercept-density},
Proposition~\ref{prop:averaged-root-hitting} proves the exact general
per-law, per-interval inequality and the bound on $C_{\mathcal D}$.

Under the separately scoped
Assumption~\ref{assump:random-intercept-witness},
Proposition~\ref{prop:random-intercept-envelope} proves cube support, the
jointly measurable conditional density, the exact reciprocal-width cap,
and the averaged envelope.  Proposition~\ref{prop:random-intercept-root-hitting}
proves the corresponding root-hitting bounds, and
Proposition~\ref{prop:fixed-width} proves every fixed-width statement.

For the explicit witness,
Propositions~\ref{prop:exact-sheet-support}
and~\ref{prop:cube-positive-width} prove exact support, cube membership,
and almost-sure positivity.  Lemma~\ref{lem:latent-recovery} and
Proposition~\ref{prop:witness-conditional-cap} prove the actual
conditioning identity, the measurable fallback construction, and the
exact nonzero-fiber cap.  Lemmas~\ref{lem:positive-mass-cap}
and~\ref{lem:exact-witness-average} prove, respectively, infinite
essential cap and the exact finite cap integral, while
Proposition~\ref{prop:witness-root-hitting} proves the singleton interval
and $C_{\{\mu_q\}}$ bounds.

Propositions~\ref{prop:sheet-affine-hull}
and~\ref{prop:sheet-null-volume} prove the exact affine dimension and
within-hull nullity.  Lemma~\ref{lem:monic-embedding} preserves those
properties in descending fixed-monic coordinates, and
Propositions~\ref{prop:affine-pushforward}
and~\ref{prop:affine-nonmembership} prove the precisely limited
full-column-rank affine-latent exclusion.  Proposition~\ref{prop:complete-sufficient-theorem}
collects these conclusions with their original quantifiers and records
the sufficient-only boundary.  Every displayed dependence in
Theorem~\ref{thm:main} is therefore proved with no hidden constant,
probability-mode conversion, or horizon change.
\end{proof}

\begin{proof}[Proof of Corollary~\ref{cor:joint-density-baseline}]
Lemma~\ref{lem:joint-density-bridge} constructs the conditional-density
versions and proves
$\bar\kappa_{\mathcal A}\leq
\kappa_{\mathrm{joint}}(2R)^{d-1}$ under the corollary's optional premise.
Proposition~\ref{prop:joint-density-baseline} then gives the interval and
class bounds and evaluates the derivative sum exactly when $B=1$,
including the $d=1$ empty-tuple case.
\end{proof}
